
Weight space learning treats neural network parameters as data, enabling analysis and manipulation of model populations. Neural Functional Networks (NFN) demonstrated that permutation-equivariant encoders can learn from model collections when all models share the same architecture, achieving performance improvements on implicit neural representation tasks (Zhou et al., 2024). Model merging techniques combine models through weight-space arithmetic but require identical architectures (Wortsman et al., 2022; Ilharco et al., 2022). The question of whether these methods extend to heterogeneous model populations spanning different architecture families remains unexplored.

The challenge arises from architecture-specific coordinate conventions. Neural networks exhibit permutation symmetries—neurons can be reordered without changing function (Hecht-Nielsen, 1990; Frankle \& Carbin, 2019)—but these symmetries differ across architectures. CNNs have spatial locality constraints, Transformers have attention head permutations, and RNNs have temporal unrolling patterns. A weight space encoding approach might project weights from different architectures into a shared space where these coordinate differences are abstracted away.

No prior work has systematically tested whether such cross-architecture weight space encodings can be learned. Git Re-Basin aligns models within single architectures using explicit permutation search (Ainsworth et al., 2022) but has not been extended to heterogeneous populations. NFN's success on homogeneous populations provides no guidance for the heterogeneous case.

We hypothesized that extending NFN's Deep Sets architecture with architecture embeddings and an explicit equivariance loss would enable cross-architecture learning. Architecture embeddings would provide family-specific context, while an MSE-based equivariance loss would encourage the encoder to factor out permutation symmetries. We tested this on synthetic model zoos with three gate metrics: reconstruction error (encoding capacity), frozen-K generalization (architecture independence), and kernel robustness (permutation invariance).

The approach failed completely. Kernel robustness reached 0.00\% (target ≥90\%), indicating that the MSE equivariance loss did not learn permutation invariance. Reconstruction error reached 19.18\% (target <10\%), indicating K=32 dimensions are insufficient. Frozen-K generalization reached 10.31\% (target <10\%), marginally above target, which combined with visualization evidence suggests architecture embeddings may prevent cross-architecture abstraction. This systematic failure across all three metrics reveals specific obstacles in the tested configuration.

We identify three root causes. First, MSE-based equivariance loss is insufficient for learning group structure—gradient descent on reconstruction objectives does not enforce the homomorphism constraint needed. Second, architecture embeddings may anchor representations to family-specific coordinates, preventing the abstraction necessary for cross-architecture transfer. Third, quotient space dimensionality requirements are underestimated—our results suggest substantially higher dimensions may be necessary, though exact requirements remain unknown.

We make the following contributions: (1) Systematic evaluation of cross-architecture weight space encoding using Deep Sets with architecture embeddings and MSE equivariance loss, establishing gate metrics that reveal failure modes. (2) Root cause analysis identifying three specific issues in the tested configuration: equivariance loss inadequacy, architecture embeddings potentially harmful, and dimensionality underestimation. (3) Alternative directions grounded in failure analysis: contrastive learning for stronger equivariance signal, architecture-agnostic encoding, increased quotient dimensions, and validation on homogeneous populations first.
